\chapter{Proposed: Wild-4D SLAM: Learning to Fuse Uncertainty and Motion Prior for Robust 4D SLAM in the Wild}
\label{chapter:sparsevio}
\fancyhead[LO]{\makebox[\textwidth][r]{An Open-Source Scene Representation Framework Towards
Long-term Autonomy}}

\begin{framed} 
Wild-4D SLAM aims to build a general visual SLAM system that sees and survives where traditional vision fails — through darkness, dynamics, degradation — unlocking reliable autonomy in the wild.
\end{framed}

\section{Introduction}

Robust 3D scene reconstruction under real-world conditions remains a grand challenge for computer vision and robotics. In ideal settings with static scenes and good lighting, purely visual SLAM (Simultaneous Localization and Mapping) systems can produce impressive 3D maps. However, no purely visual system to date can reliably handle highly \textbf{dynamic scenes or complete darkness}. Conventional visual SLAM algorithms assume mostly static environments with sufficient illumination and texture; when these assumptions break – e.g. moving objects dominate the view or the camera operates in low-light – visual methods often fail by losing track or drifting~\cite{li2024megasam}.
% This gap in robustness motivates exploring additional sensing modalities: an Inertial Measurement Unit (IMU) offers proprioceptive motion cues that can complement vision. An IMU provides accelerations and angular velocities at high frequency, which are insensitive to lighting and can serve as a short-term motion prior when visual tracking falters. In other words, \textbf{IMU motion prior can dramatically improve SLAM performance in dynamic, dark, and degraded visual conditions} by stabilizing pose estimation and guiding the vision-based mapping.


Simultaneously, the field is witnessing the emergence of “3D foundation models” – large transformer-based networks that learn powerful representations for diverse 3D vision tasks. Recent examples like DUSt3R~\cite{wang2024dust3r} and CUT3R~\cite{wang2025continuous} VGGT~\cite{wang2025vggt} demonstrate unprecedented ability to unify multiple 3D tasks (depth estimation, pose estimation, camera parameter) within one model~\cite{wang2025vggt}. These visual-geometric transformer models leverage large-scale learning to handle complexity that classic geometry methods struggle with. For instance, CUT3R can ingest a stream of RGB images and incrementally produce a consistent 4D reconstruction without needing offline optimization. Such models have achieved state-of-the-art accuracy on standard benchmarks, hinting at a new paradigm for scene reconstruction “in the wild”. Yet, notably, these foundation models today rely purely on vision – they lack integration of inertial cues and thus may still break under the extreme conditions described above (fast motions causing blur, moving obstacles, or darkness where camera frames carry little information).



This proposal aims to overcome the above limitation by integrating IMU information into 3D foundation model pipelines to achieve \textbf{real-time and resilient state estimation and 4D scene reconstruction in scenarios with dynamic objects and degraded visuals}. We posit that conditioning modern 3D transformer models with IMU data will significantly enhance their robustness, combining the best of both worlds: the learned spatial perception of foundation models with the reliability of inertial odometry. 



% Furthermore, we will develop a novel measurement-to-pose uncertainty estimation framework to adaptively fuse visual and inertial cues. By quantifying the confidence in each modality’s contribution (e.g. low under poor visual conditions, higher under stable tracking), the system can intelligently adjust how much it “trusts” the IMU vs. the camera at any given time. In summary, our goal is a new visual-inertial 4D reconstruction method that remains accurate and stable where prior vision-only methods would fail – enabling autonomous agents to map dynamic, dark environments with unprecedented reliability.

\section{Related Work}
\subsection{Visual SLAM and VIO}
Visual SLAM has a rich history, with feature-based methods like ORB-SLAM~\cite{mur2015orb} and direct methods like LSD-SLAM~\cite{engel2014lsd} and DSO~\cite{engel2017direct} enabling real-time monocular SLAM under favorable conditions. ORB-SLAM in particular demonstrated robust keypoint-based mapping and inspired many extensions including stereo, RGB-D, and inertial versions (ORB-SLAM3~\cite{orbslam3}). Visual-inertial SLAM (VI-SLAM) and visual-inertial odometry (VIO) combine camera and IMU data to improve robustness and scale observability. Traditional VIO systems such as OpenVINS\cite{openvins}, OKVIS~\cite{leutenegger2022okvis2}, and VINS-Mono/VINS-Fusion~\cite{vins-mono,qin2019general} perform sensor fusion via filter or factor graph, and are considered well-established solutions for state estimation. For instance, ORB-SLAM3 and OKVIS2 have shown accurate performance by tightly coupling IMU measurements with visual feature-based mapping. However, these methods still assume mostly static scenes and sufficient lighting (e.g. a camera blinded by darkness or many moving objects). On the dynamic scene front, approaches like DynaSLAM\cite{bescos2018dynaslam} explicitly detect and mask out moving objects (using multi-view geometry or deep learning) on top of ORB-SLAM2. By enforcing a static-map assumption through masking and even background inpainting, DynaSLAM was able to significantly outperform standard SLAM baselines in highly dynamic scenarios. Nonetheless, DynaSLAM and similar systems did not incorporate inertial cues – they focus on vision – and their runtime can be heavy due to object detection networks. Moreover, they treat the problem as \textbf{eliminating dynamics rather than modeling them}, so reconstructing a 4D (time-varying) scene is beyond their scope.



\subsection{Learning-Based 3D Foundation Models.}
Inspired by the success of large foundation models in NLP and 2D vision, researchers have started developing foundation models for 3D vision. DUSt3R is a pioneering example that introduced a transformer-based architecture capable of dense 3D reconstruction from arbitrary collections of images without known camera parameters\cite{wang2024dust3r}. Its successors have expanded these capabilities: MASt3R builds on DUSt3R to additionally perform image matching and structure-from-motion, achieving metric scale reconstruction by learning to infer scale and match keypoints in 3D~\cite{leroy2024grounding}. 
Meanwhile, CUT3R~\cite{wang2025continuous} and related models have moved toward online 4D reconstruction, maintaining a persistent state that updates with each new frame. These models represent a paradigm shift: instead of hand-engineered feature matching and incremental bundle adjustment, a learned model directly outputs structure and motion. Impressive results have been reported – for example, DUSt3R and MASt3R achieve competitive or superior accuracy on depth and pose benchmarks compared to classical SfM pipelines. However, these works are difficult to achieve robust performance in \textbf{ highly dynamic and perceptually degraded environments}. 


\begin{figure}
    \centering
    \includegraphics[width=0.5\linewidth]{figs/ch9/relatedwork.png}
    \caption{Venn diagram highlighting the research landscape. Traditional visual SLAM and classical VIO (visual-inertial odometry) address complementary parts of the problem, and learning-based 3D vision models provide new capabilities. Prior works exist in the pairwise intersections (e.g. visual-inertial SLAM, learned visual SLAM methods, and learned VIO, but none covers the intersection of all three. Wild-4D SLAM targets this unexplored triple intersection, aiming to integrate visual SLAM, IMU-based prior, and 3D foundation model learning into a single resilient system.}
    \label{fig:enter-label}
\end{figure}

\subsection{Resilience in Perception (Dynamic and Low-Visibility Conditions)}

Achieving robustness in dynamic environments has been a focus of many recent SLAM improvements. One strategy is to integrate model-based motion prior or outlier rejection schemes. For instance, DynaSLAM~\cite{bescos2018dynaslam} and other dynamic SLAM frameworks remove feature correspondences that violate rigid scene assumptions. More recently, deep visual SLAM frameworks like DROID-SLAM~\cite{teed2021droid} and MegaSaM~\cite{li2024megasam} incorporate learned components to handle moderate scene dynamics. MegaSaM in particular extends a differentiable bundle-adjustment approach (DROID-SLAM) by integrating learned monocular depth priors and motion segmentation, allowing it to track cameras through videos with moving objects and even low-parallax motions that would break traditional SfM. Despite these advances, fully solving dynamic scenes with vision alone remains hard – there are always cases (e.g. most of the view is moving and darkness).

\textbf{Summary:} There is a clear opportunity at the intersection of 3D foundation models and visual-inertial SLAM. On one hand, transformer-based 3D models bring powerful learned priors for geometry and can handle aspects like uncalibrated cameras and even moderate dynamics in a purely data-driven way. On the other hand, IMUs offer strong physical cues and failure-case coverage (darkness, rapid motion) that vision alone cannot handle. Yet, so far these two trends have not been fully combined – current 3D foundation models do not incorporate inertial data, and conversely, classical VI-SLAM has not benefited from the advances in foundation models. This research \textbf{aims to fill that gap by designing an IMU-conditioned 3D reconstruction framework that inherits the strengths of both. We aim to show that foundation-model priors can yield a resilient, computationally efficient 4D reconstruction system.}


\section{Contribution}

We introduce Wild-4D SLAM. Our key contribution is listed as follows:
\begin{itemize}
    \item \textbf{IMU-dominated SLAM Pipeline:}
    A learning-based approach that treats inertial sensing as the primary sensor of state propagation, while incorporating sparse visual updates from a 3D foundation model to correct global drift.
    \item \textbf{IMU-conditioned Visual Ground Transformer:} the first IMU-conditioned Visual Geometry Transformer (IVGT) that tightly integrates inertial motion priors into the visual-geometric feature. It provides global motion awareness to local visual features, resulting in more resilient and generalizable feature representations.
    \item \textbf{Learning based Adaptive Sensor Fusion:} a hybrid approach that combines analytic uncertainty estimation with data-driven uncertainty prediction on pose uncertainty which is less explored.
    \item \textbf{Significant Improvement on Robustness and Efficiency:} Our solution can potentially overcome darkness, dynamic and other extreme degradation environments that have not been achieved before and maintain high efficiency.
\end{itemize}








% Our investigation is guided by four research questions:

% \begin{itemize}
%     \item \textbf{Foundation-Model Integration} – How can a large-scale, pretrained foundation model be integrated into a SLAM system to provide strong perceptual priors?
%     \item \textbf{Resilience and Recovery} – How can this integration endow a VIO framework with built-in recovery mechanisms that handle sensor dropouts and dynamic scene changes?
%     \item \textbf{Efficiency vs. Robustness} – How can we leverage foundation models in a way that preserves—or improves—runtime efficiency while boosting robustness?
%     \item \textbf{Performance in Degraded Conditions} – How can foundation-driven priors help visual odometry overcome its traditional failure modes in perceptually degraded environments (e.g., low light, motion blur, obscurants)?
%     \item \textbf{Uncertainty Estimation on poses}- How to propagate the uncertainty from the measurements to poses?  
% \end{itemize}






% Accurate and robust tracking is a fundamental requirement for virtual and augmented reality (VR/AR) systems, where even small localization errors can degrade user experience by causing visual odometry drift.

% We will begin our experiments with VR/AR systems, which demand continuous, robust, and precise visual-inertial odometry (VIO). Traditional VIO methods depend heavily on visual features, leaving them vulnerable in low-light settings, textureless scenes, and during rapid motion. In contrast, IMU-dominated SLAM exploits the high-frequency motion estimates from inertial sensors, reducing reliance on visual cues while preserving localization accuracy.

% To this end, we present Sparse-VIO, an IMU-centric SLAM framework tailored for reliable tracking in VR/AR applications. Rather than continuously fusing camera and IMU data, Sparse-VIO prioritizes inertial propagation and invokes visual relocalization only intermittently to correct long-term drift. This design is especially advantageous for head-mounted and mobile VR/AR devices, where power efficiency and resilience to visual degradation are critical. An adaptive fusion mechanism modulates the contribution of vision and inertial data according to environmental conditions, delivering consistent tracking performance even in challenging scenes. Our approach makes four key contribution:

% \begin{itemize}
%     \item \textbf{IMU-Centric Sensor Fusion Backbone}:
% Wild-4D SLAM elevates the IMU to the primary source of motion propagation, providing resilience to perceptual degradation and ensuring reliable state estimation in challenging environments.

% \item \textbf{Sparse Visual Corrections via Foundation Model Priors}:
% Instead of dense visual tracking, the system leverages sparse visual updates informed by a 3D foundation model to perform global pose correction, significantly reducing visual processing demands while preserving accuracy.

% \item \textbf{Resilience in Degraded Conditions}:
% By relying on robust IMU-only odometry during sensor dropouts or degraded perception, Wild-4D SLAM maintains continuous operation and avoids failure where traditional visual-inertial systems would collapse.

% \item \textbf{Online Uncertainty Estimation}:
% The design minimizes computational overhead by avoiding continuous dense vision processing, making it well-suited for real-time deployment on resource-constrained platforms.
% \end{itemize}

% \begin{itemize}
%  \item \textbf{Infrequent Visual Relocalization} – 
%  Instead of continuously tracking features for all the images,  we only 
%  use sparse visual images to correct the long-term IMU odometry drift.  This sparse image view will significantly reduce computational load while preserving global consistency.  

% \item \textbf{IMU-centric Sensor Fusion} – The system dynamically selects motion priors from either vision or IMU, ensuring robustness in different kind of visual challenges such as aggressive motion, low-light conditions, and feature-sparse environments—addressing key challenges in VR/AR applications.  

% \item \textbf{Dense Mapping with Sparse Views} – We leverage 3D foundation model like Mast3R\cite{leroy2024grounding} or VGGT\cite{wang2025vggt} to achieve dense reconstruction of environments with varying numbers of sparse views. Our approach estimates accurate camera parameters and generates dense point clouds of the environment.


% \item \textbf{Overcome Very Extreme Environments} – Since our pipeline is using the IMU foundation model, it remains unaffected by complete darkness and highly dynamic environments.  The sensor fusion pipeline will directly bypass the visual sensor solution and adopt a solution from the IMU foundation model in extreme environments.
     
    
% \end{itemize}

\section{Proposed Method}


 \begin{figure*}[h!]
\setlength{\abovecaptionskip}{0.1cm}\centering
\includegraphics[width=\textwidth]{figs/ch9/system.png} 
\caption{\textbf{The pipeline of Wild-4D SLAM.} We propose:  
(\textbf{a}) \textbf{\textit{IMU Foundation Model}}: Building on our existing work, TartanIMU, we adapt the IMU model to human motion pattern and output high-frequency relative poses.  
(\textbf{b}) \textbf{\textit{Visual Foundation Model}}: Leveraging Mast3R, we enable relocalization with extremely sparse observations, correcting drift from the IMU Foundation Model.  
(\textbf{c}) \textbf{\textit{IMU-centric Sensor Fusion}}: Our IMU-centric sensor fusion pipeline dynamically adjusts reliance on vision-based priors or IMU-based priors and achieves robust sensor fusion according to the change of environments.}  
	\label{fig:teaser}\vspace{-2mm}
\end{figure*}

% \begin{itemize}
% \item \textbf{IMU Foundation Model:} We introduce the IMU foundation model, which features a shared backbone designed for scalable, cross-platform motion estimation, as illustrated in \ref{fig:teaser}. The model architecture comprises four main components: \textit{data alignment and augmentation, a heterogeneous shared IMU backbone, a multi-head structure for different robots, and a specialized loss function}.  
% Our proposed IMU foundation model builds upon the previous work TartanIMU and will be trained on a diverse human motion dataset from VR/AR systems.  

% \item \textbf{Visual Foundation Model:} We present a real-time visual odometry system based on the two-view 3D reconstruction prior such as MASt3R \cite{leroy2024grounding} or VGGT\cite{wang2025vggt} as shown in \fref{fig:teaser}. The system operates without assumptions about the camera model for each image and can perform tracking, mapping, and relocalization, even in scenarios with minimal overlap between views. This capability enables infrequent yet effective localization.  

% \item \textbf{IMU-centric Factor Graph:} The design of our system is driven by a key insight: an inertial measurement unit (IMU) provides smooth, noise-contaminated measurements with minimal outliers. This makes IMU-based estimates highly accurate, provided that bias drift is well-constrained by other sensors. We leverage our Super Odometry sensor fusion pipeline \cite{zhao2024subt}, ensuring that as long as the visual foundation models and IMU foundation models supply reliable relative pose priors to constrain the IMU pre-integration factor, all sensor measurements can be successfully fused.
% \end{itemize}


\subsection{IMU-Conditioned Visual Geometry Transformer Architecture:}
We propose the first IMU-conditioned Visual Geometry Transformer (IVGT) that tightly integrates inertial motion priors into the visual-geometric feature reasoning process. Existing VGGT (such as Dust3R/Cut3R) architectures rely solely on visual inputs and 3D geometry, making them vulnerable to degradation in dynamic, dark, or perceptually challenging environments where visual features become unreliable. 

This contribution includes developing a method to encode IMU measurements and inject them into the model’s latent representation of the scene. For example, we may introduce an IMU encoder that processes the sequence of accelerometer/gyroscope data between two frames and produces a \textbf{pose prior that conditions the transformer’s self-attention layers or the pose estimation head}. 

The IMU-conditioned VGGT architecture provides global motion awareness to local visual features, resulting in more resilient and generalizable feature representations. We explore several strategies for this integration:

\begin{figure}
    \centering
    \includegraphics[width=1.0\linewidth]{figs/ch9/vggt.png}
    \caption{Two ways incorporate the IMU latent features into Visual Geometry Transformer Architectures}
    \label{fig:vggt}
\end{figure}

\paragraph{Latent State Injection (Cross-Modal Attention)}
We inject IMU latent features directly into the token sequence processed by the transformer. The IMU latent vector $\mathbf{z}_{\text{imu}}$ is embedded as a learnable token $\mathbf{t}_{\text{imu}}$, concatenated to the visual and geometric tokens shown as\fref{fig:vggt}:

\begin{equation}
\mathcal{T} = \{ \mathbf{t}_{1}, \mathbf{t}_{2}, \ldots, \mathbf{t}_{N}, \mathbf{t}_{\text{imu}} \}
\end{equation}

where $\mathcal{T}$ is the token set, and $N$ is the number of visual and geometric tokens. The transformer layers then allow cross-attention between visual, geometric, and IMU tokens, enabling dynamic adaptation of feature aggregation based on motion context.

\subsection{Feature Modulation (FiLM-style Conditioning)}

In this strategy, the IMU latent features are used to modulate the visual features via Feature-wise Linear Modulation (FiLM) shown as \fref{fig:vggt}. Given a visual feature map $\mathbf{F}_{\text{vis}}$ and an IMU latent code $\mathbf{z}_{\text{imu}}$, we generate scaling and bias parameters:

\begin{equation}
(\gamma, \beta) = \text{MLP}(\mathbf{z}_{\text{imu}})
\end{equation}

The modulated feature is then:

\begin{equation}
\widetilde{\mathbf{F}}_{\text{vis}} = \gamma \odot \mathbf{F}_{\text{vis}} + \beta
\end{equation}

where $\odot$ denotes element-wise multiplication. This conditioning can be applied before each transformer block or selectively to attention or feedforward layers.

\paragraph{Summary}
Each of these strategies offers a different trade-off between model complexity and effectiveness:
\begin{itemize}
    \item \textbf{Feature Modulation:} Lightweight and simple, effective for small perturbations.
    \item \textbf{Latent State Injection:} Richer interaction, especially for strong dynamic changes.
\end{itemize}




% We will explore architectural innovations like inertial-conditioned attention (where IMU-derived motion estimates influence the attention weights for matching features across frames) and inertial pose warping (using IMU integration to initially align point cloud predictions before refinement by the network). 



% The resulting architecture will be the first 3D foundation model that \textbf{natively fuses vision and inertial data}, maintaining the generality of foundation models (e.g. working with arbitrary uncalibrated image sequences) while gaining physical robustness. 


% We anticipate this will dramatically improve performance in challenging settings, as the IMU provides strong cues to camera motion when visual information is weak. By leveraging the IMU, the model can, for instance, keep tracking through a dark interval and quickly re-lock onto visual features when they reappear, thereby preventing catastrophic drift or failure.


\begin{figure}
    \centering
    \includegraphics[width=1.0\linewidth]{figs/ch9/pose_uncertainty.png}
    \caption{The proposed pose covariance architecture, we predict the lower-diagonal matrix and perform the inverse Cholesky decomposition to ensure the positive semi-definite constraint of the
covariance matrix. The predicted covariance is merged with the output of the covariance model for the potential function of factor graphs.}
    \label{fig:pose_uncertainty}
\end{figure}


\subsection{Learning based Adaptive Backend Optimization (From measurement Uncertainty to Pose Uncertainty):}

Mathematically, uncertainty is typically modeled as a Gaussian distribution; however, this assumption may break down in degraded or highly dynamic environments, where real-world noise does not necessarily follow a Gaussian form. In neural networks, uncertainty estimation is prone to overfitting, particularly when the dataset does not cover the full spectrum of environmental variability.

Nevertheless, accurate uncertainty estimation is crucial for achieving reliable sensor fusion and state estimation. Existing pipelines primarily focus on estimating uncertainty at the measurement level (e.g., feature tracks or depth maps). However, relatively few works address uncertainty at the pose level, despite its critical importance for efficient and robust sensor fusion.

In particular, if pose uncertainty is available during fusion, it enables significantly more effective weighting of visual and inertial modalities and reduces computational demands on backend optimization (e.g., fewer iterations or simpler marginalization strategies).

To address this gap, we propose a hybrid approach that combines analytic uncertainty estimation with data-driven uncertainty prediction. To avoid overfitting the uncertainty prediction to specific datasets, we adopt the following strategies:

\begin{itemize}
    \item \textbf{Learning Pose Uncertainty from Measurement Uncertainty}: Rather than directly supervising pose uncertainty using only pose error, we propose to learn pose uncertainty as a function of measurement uncertainties (e.g., depth uncertainty, optical flow uncertainty) produced by the VGGT network. Specifically, we will develop a learned encoder-decoder architecture that maps measurement uncertainties to pose uncertainties. This provides a richer learning signal and better captures how local measurement noise aggregates into global pose errors.
    \item \textbf{Integration with Analytic Uncertainty}:  As illustrated in Fig.~\ref{fig:pose_uncertainty}, the final pose uncertainty will be modeled as a combination of analytic uncertainty (derived from classical uncertainty propagation) and data-driven uncertainty (learned by the network). This hybrid formulation is expected to stabilize training and produce better-calibrated uncertainty estimates.
    \item \textbf{Differentiable IMU-Centric Factor Graph Supervision}: We will use an IMU-centric factor graph optimization layer, designed to be differentiable, to backpropagate pose estimation errors and uncertainty calibration errors. This supervision will tune both the predicted pose uncertainty and the intermediate measurement uncertainties.
\end{itemize}

While the above discussion focuses primarily on visual uncertainty, we will also extend the framework to incorporate IMU uncertainty. Through factor graph optimization, the network will dynamically adjust both visual and inertial uncertainty estimates, enabling more intelligent fusion at the pose level.

We believe that this learning-based adaptive fusion mechanism will significantly improve the system's resilience. For example, if the robot enters a dark tunnel, the model will recognize the increased visual uncertainty and rely more heavily on inertial measurements; when strong visual cues return, the system can decrease reliance on the IMU and correct accumulated drift through visual loop closures.

By the end of the project, we aim to demonstrate that this uncertainty-driven fusion strategy consistently outperforms naive fusion methods (e.g., fixed trust in sensors) across a range of dynamic and degraded scenarios, providing a generalizable and robust solution for multi-sensor integration in learned SLAM systems.


% Benchmarking and Empirical Analysis of Resilience: (In addition to the two main theoretical contributions above, an important outcome will be a thorough evaluation and new insights into failure modes and performance of visual-inertial SLAM in extreme conditions. This is a more practical contribution.) We will contribute an extensive benchmarking of our proposed system against prior methods on challenging datasets, analyzing where traditional visual SLAM breaks and how much our IMU-conditioned model improves. By quantitatively measuring performance in dynamic low-light scenarios, we aim to highlight the specific gains from IMU integration. This empirical study will serve as a reference for the community on the benefits of multi-modal fusion in next-generation SLAM systems. Moreover, we will release open-source code and guidelines for integrating IMUs into learned 3D reconstruction models, lowering the barrier for future researchers to build upon our approach.



\section{Benchmark Plan}

To validate the performance of Wild-4D SLAM, we will carry out an extensive evaluation on both public benchmark datasets and real-world collected data. Our benchmarking plan covers datasets, baseline methods for comparison, and evaluation metrics as follows:

\subsection{Datasets:} We will use a mix of standard and challenging datasets that together cover the “wild” scenarios of interest:

\textbf{EuRoC MAV Visual-Inertial Dataset\cite{burri2016euroc}}: a popular VIO benchmark captured by a drone in indoor environments (Vicon Room, industrial machine hall) with fast motion and varying illumination. EuRoC provides stereo images, IMU readings, and ground-truth motion, enabling evaluation of accuracy on trajectories. Although largely static, EuRoC sequences include difficult conditions like texture-poor areas and aggressive motion where our IMU integration and adaptive fusion will be stress-tested.

\textbf{LaMAR Dataset \cite{sarlin2022lamar}:} A benchmark with a comprehensive capture (lidar,imu, camera) and GT pipeline that co-registers realistic
trajectories and sensor streams captured by heterogeneous AR devices in large, unconstrained scenes.

\textbf{VBR Benchmark\cite{brizi2024vbr}:} an outdoor lidar-visual-inertial dataset with handheld motion, including sequences with dynamic people and HDR lighting changes. This will test our system’s dynamic scene handling and robustness to exposure variations. We will specifically use sequences where a person walks in front of the camera or where lighting is switched on/off, to evaluate dynamic mapping capability.

% KITTI Odometry (with IMU): outdoor driving sequences with a forward-looking stereo camera and IMU data (available through the KITTI raw data). KITTI provides dynamic real-world scenes (moving cars, pedestrians) and various lighting (daytime, shadows). This lets us evaluate large-scale outdoor performance and compare to state-of-the-art VIO on driving data. Subsets recorded at dawn/dusk will serve as mild low-light tests.


\textbf{Custom “Wild” Sequences:} We will collect or additional sequences to specifically evaluate extreme cases: (i) Low-light environment – e.g., a sequence in an underground garage at night with only headlights or a few lamps, testing SLAM in near-darkness. (ii) Smoke/Fog conditions – using either a fog machine in a lab environment or a synthetic data simulation with fog, to mimic sensor-degraded conditions (here thermal would normally be used, but our method should cope to some extent with just camera+IMU). (iii) Highly dynamic scene – e.g., a sequence in a populated hallway where dozens of people move, or a street scene with heavy traffic, to test the limits of dynamic object handling. These custom sequences, although not having precise ground truth, will be useful for qualitative evaluation and stress-testing robustness (we may use approximate truth from motion-capture for indoor or GPS for outdoor for quantitative metrics).


\subsection{Baselines:} We will compare Wild-4D SLAM against several representative baseline methods:

\begin{itemize}
    \item \textbf{Classic VIO Method: } ORB-SLAM3 DynaSLAM 
    \item \textbf{Learning Methods:} DROID-SLAM, MegaSAM, Mac-VO
    \item \textbf{3D Foundation model: }VGGT, Dust3R, Cut3R
    \item  \textbf{3D Foundation model SLAM:} Mast3R-SLAM
\end{itemize}


\subsection{Metric:}
\paragraph{Localization Accuracy:} The primary metric will be Absolute Trajectory Error (ATE) – the root-mean-square error between the estimated trajectory and ground truth, after alignment. We will compute ATE for each sequence (in meters for translation, degrees for rotation) and also report Relative Pose Error (RPE) over short intervals to capture drift per unit distance.

\paragraph{Mapping Quality:} For datasets where ground-truth maps or point clouds are available (e.g., KITTI has LiDAR maps, or for EuRoC we can use the dense multi-view stereo reconstruction as pseudo-ground truth), we will evaluate the quality of reconstruction. Metrics like the average distance of mapped points to ground-truth structure (ATE of map points) or intersection-over-union (IoU) of occupied space in a volumetric map can be used. 
 
\paragraph{Robustness Metric\cite{zhao2024subt}:} We will record the number of tracking failures or resets needed for each system. A robust SLAM should ideally run through a whole sequence without losing track. We will compare the fraction of sequences successfully completed. Furthermore, we plan to quantify uncertainty calibration: e.g., checking if our uncertainty estimates correlate with actual errors (an ideally uncertainty-aware system will give higher uncertainty before failures). This can be done by plotting uncertainty vs. localization error and computing correlation.
Computational 

\paragraph{Performance:} We will measure runtime (frame processing time, CPU/GPU usage) and memory consumption. We need to ensure the system runs close to real-time. We will report average frame rate achieved on a standard computer or onboard computer, and compare to baseline runtimes. 




By the end of this evaluation campaign, we expect to clearly demonstrate that Wild-4D SLAM outperforms traditional visual SLAM, classical VIO, and existing learning-augmented SLAM methods in the targeted challenging conditions, thus validating our research hypothesis.












% In addition to the proposed IMU-centric fusion design, this work extends the theoretical foundation of Wild-4D SLAM by developing a principled method for estimating pose uncertainty from measurement-level uncertainties. Specifically, we aim to propagate uncertainty from depth flow and patch-level optical flow into the overall pose estimation process shown in \fref{fig:pose_uncertainty}. Unlike traditional VIO systems that often treat visual updates deterministically, our framework explicitly models the confidence of each sparse visual constraint. This allows the system to dynamically weigh each correction based on its reliability, improving robustness in regions with weak texture, ambiguous depth, or degraded lighting.



\section{Future Work Discussion}
  While our current approach implicitly handles dynamics via sensor cues, future work could explicitly integrate semantic understanding. For example, a learned model could be extended to output not just geometry but also labels or instance segments for dynamic objects (people, cars) and treat them separately in the map. This would enable 4D dynamic scene understanding – mapping not only the static world but also tracking trajectories of moving agents. Incorporating an IMU in such a setting could help estimate the ego-motion and segregate it from object motions more cleanly. High-level reasoning (e.g., recognizing an object and knowing it can move) could further improve robustness and could be an extension after the core sensor fusion is in place.





\section{Proposed Timeline for Projects}

\begin{itemize}
    \item \textbf{07/15:} Finalize baseline methods and benchmarks for evaluation.
    \item \textbf{08/20:} Complete evaluation of baseline methods on public datasets.
    \item \textbf{08/25:} Train the IMU foundation model with a focus on human motion patterns.
    \item \textbf{08/30:} Collect datasets in dynamic and low-visibility (e.g., darkness) environments.
    \item \textbf{09/05:} Integrate the IMU foundation model with a 3D visual foundation model.
    \item \textbf{09/30:} Investigate and implement pose uncertainty estimation techniques.
    \item \textbf{10/15:} Integrate the uncertainty estimation module into the Wild-4D SLAM framework.
    \item \textbf{10/20:} Complete the first full draft of the paper.
    \item \textbf{10/25:} Revise and finalize the second draft based on internal feedback.
    \item \textbf{11/05:} Begin producing the submission video and writing supplementary materials.
    \item \textbf{11/15:} Finalize and submit all materials for CVPR 2026.
\end{itemize}